\section{Theory: Safety and Stability Guarantees}

Let $s_t$ be the observed state, $a_t^{\pi}$ the learned policy action, $a_t^{h}$ the fallback heuristic action, and $u_q(s_t)$ the quantum uncertainty estimate. Let the switching rule be
\[
\hat a_t=
\begin{cases}
a_t^{h}, & u_q(s_t)>\tau_q,\\
a_t^{\pi}, & u_q(s_t)\le\tau_q.
\end{cases}
\]
Let $\mathcal{S}(s,a)$ be the shield operator that enforces hard constraints (bounds and rate limits) and outputs an executable action $a_t=\mathcal{S}(s_t,\hat a_t)$.

\paragraph{Theorem 1 (Safety Bound with Shielded Quantum Gating).}
Assume the shield is complete for hard constraints, i.e., for every state-action pair $(s,a)$ the output $\mathcal{S}(s,a)$ satisfies all hard safety constraints. Then:
\begin{enumerate}
	\item Within the simulation environment, and assuming a bug-free and complete shield implementation, the probability of catastrophic action is zero.\footnote{This is an in-simulator guarantee, not a real-world absolute guarantee. Deployment can fail due to software defects, model mismatch, sensing faults, or incomplete constraint encoding; therefore operational risk controls remain necessary.}
	\item Let $Q_t(\cdot)$ be one-step expected utility and let $a_t^{\star}=\arg\max\{Q_t(a_t^{\pi}),Q_t(a_t^h)\}$. If quantum uncertainty calibration error is bounded by $\epsilon_{\text{quantum}}$ and the shield induces utility deviation at most $\epsilon_{\text{shield}}$, then
	\[
	\mathbb{E}[Q_t(a_t^{\star})-Q_t(a_t)] \le \epsilon_{\text{quantum}}+\epsilon_{\text{shield}}.
	\]
\end{enumerate}

\textit{Proof.}
For part (1), define the hard-constraint-satisfying action set in state $s$ as $\mathcal{A}_{\mathrm{safe}}(s)$. Completeness of $\mathcal{S}$ states
\[
\forall (s,a),\quad \mathcal{S}(s,a)\in \mathcal{A}_{\mathrm{safe}}(s).
\]
At time $t$, executed action is $a_t=\mathcal{S}(s_t,\hat a_t)$, hence $a_t\in\mathcal{A}_{\mathrm{safe}}(s_t)$ almost surely. Catastrophic actions are exactly actions outside $\mathcal{A}_{\mathrm{safe}}(s_t)$, therefore
\[
\Pr(a_t\notin\mathcal{A}_{\mathrm{safe}}(s_t))=0.
\]

For part (2), write one-step regret
\[
R_t := Q_t(a_t^{\star})-Q_t(a_t)=Q_t(a_t^{\star})-Q_t(\mathcal{S}(s_t,\hat a_t)).
\]
Add and subtract $Q_t(\hat a_t)$:
\[
R_t=\underbrace{Q_t(a_t^{\star})-Q_t(\hat a_t)}_{\text{switching regret }R_t^{\mathrm{sw}}}
+\underbrace{Q_t(\hat a_t)-Q_t(\mathcal{S}(s_t,\hat a_t))}_{\text{shield regret }R_t^{\mathrm{sh}}}.
\]
By assumption, shield projection loss satisfies $R_t^{\mathrm{sh}}\le \epsilon_{\text{shield}}$ almost surely, hence $\mathbb{E}[R_t^{\mathrm{sh}}]\le \epsilon_{\text{shield}}$.

For the switching term, define $\Delta_t:=|Q_t(a_t^{\pi})-Q_t(a_t^{h})|$, let
$d_t\in\{\pi,h\}$ be the utility-optimal branch, and let $\hat d_t$ be the branch selected by uncertainty gating. Then
\[
R_t^{\mathrm{sw}} = \Delta_t\,\mathbf{1}\{\hat d_t\neq d_t\}.
\]
By the calibration assumption for $u_q$, expected misrouting loss is bounded:
\[
\mathbb{E}[\Delta_t\,\mathbf{1}\{\hat d_t\neq d_t\}]\le\epsilon_{\text{quantum}},
\]
so $\mathbb{E}[R_t^{\mathrm{sw}}]\le\epsilon_{\text{quantum}}$.

Combining both terms,
\[
\mathbb{E}[R_t]
=\mathbb{E}[R_t^{\mathrm{sw}}]+\mathbb{E}[R_t^{\mathrm{sh}}]
\le \epsilon_{\text{quantum}}+\epsilon_{\text{shield}}.
\]
This proves part (2).

\paragraph{Corollary 1 (Graceful Degradation under Mismatch).}
Let calibration mismatch scale by factor $\kappa\ge 1$ so that effective uncertainty-routing error is $\epsilon'_{\text{quantum}}=\kappa\epsilon_{\text{quantum}}$, and let shield mismatch add residual execution error $\epsilon'_{\text{shield}}=\epsilon_{\text{shield}}+\epsilon_{\text{mismatch}}$. Then Theorem 1 implies
\[
\mathbb{E}[Q_t(a_t^{\star})-Q_t(a_t)]
\le
\kappa\epsilon_{\text{quantum}}+\epsilon_{\text{shield}}+\epsilon_{\text{mismatch}}.
\]
Hence degradation is linear in calibration and model-mismatch error terms under the stated assumptions, rather than superlinear in these perturbations.

\paragraph{Corollary 2 (Observation-Noise Robustness).}
Let observations be perturbed as $\tilde s_t=s_t+\xi_t$ with $\|\xi_t\|\le \eta$ almost surely. Assume branch-value maps are locally Lipschitz in observation space with constants $L_{\pi}$ and $L_h$, and shield projection loss remains bounded by $\epsilon_{\text{shield}}$. Then one-step regret under noisy observations satisfies
\[
\mathbb{E}\big[Q_t(a_t^{\star})-Q_t(a_t)\big]
\le
\epsilon_{\text{quantum}}+\epsilon_{\text{shield}}+(L_{\pi}+L_h)\eta.
\]
Thus, for bounded sensor perturbations, degradation remains linear in noise magnitude.

\paragraph{Theorem 2 (Variance Reduction under Risk-Penalized Return).}
Define shaped reward
\[
\tilde r_t=r_t-\lambda_{\mathrm{risk}}\,\mathrm{Var}(r_{t-k:t}), \qquad \lambda_{\mathrm{risk}}>0.
\]
Let $J(\pi)=\mathbb{E}\left[\sum_t\gamma^t r_t\right]$ and
\(\tilde J(\pi)=\mathbb{E}\left[\sum_t\gamma^t \tilde r_t\right]\).
If an update from $\pi$ to $\pi'$ satisfies $\tilde J(\pi')\ge \tilde J(\pi)$ and $J(\pi')\approx J(\pi)$, then
\[
\mathbb{E}\big[\mathrm{Var}_{\pi'}(r_{t-k:t})\big] \le
\mathbb{E}\big[\mathrm{Var}_{\pi}(r_{t-k:t})\big].
\]

\textit{Proof.}
Define discounted variance functional
\[
V(\pi):=\mathbb{E}_{\pi}\left[\sum_{t=0}^{\infty}\gamma^t\,\mathrm{Var}(r_{t-k:t})\right].
\]
By definition of shaped reward,
\[
\tilde J(\pi)=J(\pi)-\lambda_{\mathrm{risk}}V(\pi).
\]
Since $\tilde J(\pi')\ge \tilde J(\pi)$,
\[
J(\pi')-\lambda_{\mathrm{risk}}V(\pi')
\ge
J(\pi)-\lambda_{\mathrm{risk}}V(\pi).
\]
Rearrange:
\[
\lambda_{\mathrm{risk}}\big(V(\pi')-V(\pi)\big)
\le
J(\pi')-J(\pi).
\]
Because $\lambda_{\mathrm{risk}}>0$,
\[
V(\pi')-V(\pi)
\le
\frac{J(\pi')-J(\pi)}{\lambda_{\mathrm{risk}}}.
\]
Under the assumption $J(\pi')\approx J(\pi)$, let $\delta:=J(\pi')-J(\pi)$ with $|\delta|\ll 1$. Then
\[
V(\pi')-V(\pi)\le \frac{\delta}{\lambda_{\mathrm{risk}}}.
\]
In the practically relevant case $\delta\le 0$ (or $\delta$ negligible relative to $\lambda_{\mathrm{risk}}$), we obtain $V(\pi')\le V(\pi)$. Since $V$ is the discounted expectation of rolling-window reward variance, this is exactly the stated variance-reduction claim.

Equivalently, if one tracks the per-window expectation $\bar v(\pi):=\mathbb{E}[\mathrm{Var}(r_{t-k:t})]$ under stationarity or finite-horizon averaging, the same inequality carries through up to the constant discount normalization, yielding
\[
\bar v(\pi')\le \bar v(\pi).
\]
Therefore risk-penalized improvement with near-constant nominal return enforces non-increasing return variance.

These theorems justify why the combined policy (quantum-gated fallback + shield + risk shaping) improves safety behavior while preserving competitive reward in adversarial governance control.
